\documentclass{article}
\usepackage{graphicx} % Required for inserting images

\title{Online Handwriting}
\author{Hudson Mitchell-Pullman}
\date{October 2025}

\begin{document}

\maketitle

Representing Online Handwriting for Recognition in Large Vision-Language Models

Using Vision-language models to recognize online handwriting (digital ink) is more effective than images and OCR.


RQ: How can users switch between modalities and turn handwritten notes into printed text?

VLMs contain an underlying language model, so while they’re highly effective, they’re also computationally expensive.

We need a time-ordered sequence because it is more effective than OCR.
Based on images
Based on time-ordered sequences
Make the sequence of points a discrete sequence of text tokens that can be consumed by VLMs
Requires fine-tuning or parameter-efficient tuning
NO RETRAINING

We let $S$ denote a Stroke, which is a sequence of triplets $(x, y, t)$ where $t$ denotes time and $x$ and $y$ are coordinates. We let $I$ denote Ink, where $I=[s_0, s_1, ..., s_n]$. The input of the model is ink $I$, and the output is the text written in the ink $I$ as text.

Architecture:
PaLI and PaLM-E architectures.
Tokenize each digit separately. 
Found it crucial to consider time sampling, scale normalization, coordinate representation, discretization, and token dictionary.

If we visually encode, use vanilla ViT as the most popular
and standard vision encoder solution.

The performance of both VLMs fine-
tuned for handwriting recognition is comparable to or
better than state-of-the-art results.







\end{document}
